\notepage{34}{The Parity Selection Rule}
Let states and operators have definite parity:
\[
 \Pi\ket{\psi_\alpha}=(-1)^\alpha\ket{\psi_\alpha},\qquad
 \Pi^\dagger O_\beta\Pi=(-1)^\beta O_\beta,\qquad
 \alpha,\beta\in\{0,1\}.
\]
Inserting the identity twice,
\begin{align*}
 \mel{\psi_\alpha}{O_\beta}{\psi_\gamma}
 &=\mel{\psi_\alpha}{\Pi^\dagger\Pi O_\beta\Pi^\dagger\Pi}{\psi_\gamma}\\
 &=(-1)^\alpha(-1)^\beta(-1)^\gamma
       \mel{\psi_\alpha}{O_\beta}{\psi_\gamma}.
\end{align*}
If the product of the three parities is odd, then
\[
 \alpha+\beta+\gamma\ \hbox{odd}
 \quad\Longrightarrow\quad
 \mel{\psi_\alpha}{O_\beta}{\psi_\gamma}
 =-\mel{\psi_\alpha}{O_\beta}{\psi_\gamma}=0.
\]
\smallheading{Example: a Gaussian with a travelling phase}
Consider again
\[
 \psi_u(x)=\ee^{\ii p_0x/\hbar}\ee^{-x^2/(2\Delta^2)}.
\]
When $p_0=0$ the state is even. Both $X$ and $P$ are odd, so their
expectations vanish by the selection rule.
When $p_0\ne0$, the state is not a parity eigenstate, so that rule cannot
be applied directly. Nevertheless, the density is even:
\begin{align*}
 \langle X\rangle
 &=\frac1{\braket{\psi_u}{\psi_u}}
       \int_{-\infty}^{\infty}\psi_u(x)^*x\psi_u(x)\,\dd x\\
 &=\frac1{\braket{\psi_u}{\psi_u}}
       \int_{-\infty}^{\infty}x\ee^{-x^2/\Delta^2}\,\dd x=0.
\end{align*}
Equivalently, inserting $\Pi^\dagger\Pi$ on either side of $X$ changes
$x$ into $-x$ and $\psi_u(x)$ into $\psi_u(-x)$; their squared moduli
are equal. Thus the transformed integral is $-\langle X\rangle$.
For momentum, direct differentiation instead gives
\begin{align*}
 -\ii\hbar\psi_u'(x)
 &=\left(p_0+\frac{\ii\hbar x}{\Delta^2}\right)\psi_u(x),\\
 \langle P\rangle&=p_0+\frac{\ii\hbar}{\Delta^2}\langle X\rangle=p_0.
\end{align*}
\editorial{The source ends with ``similarly, $\langle P\rangle=0$.''
That conclusion is valid for $p_0=0$ only; the travelling-phase example
has $\langle P\rangle=p_0$.}
